\chapter{Neural Network Training}

\section{Simulation Environment and Implementation}

\subsection{Physics Simulation Setup}

The simulation environment is built using Tudatpy, an open-source astrodynamics
toolbox that provides high-fidelity spacecraft dynamics modeling. The physics
simulation operates in an empty universe configuration with the Solar System
Barycenter (SSB) as the reference point, eliminating gravitational
perturbations to focus purely on the spacecraft control problem. A
6-degree-of-freedom (6-DoF) model is used without mass propagation, assuming
constant mass for simplicity.

The equations of motion (EoM) are numerically integrated using a variable
step/order Bulrisch-Stoer method \autocite{bulirsch1966numerical} with fixed
tolerance and initial timestep. The integrator settings employ block-wise
integration to handle the coupled translational and rotational dynamics of the
spacecraft system. The integrator and interpolator settings were validated using Richardson's extrapolation method to ensure numerical accuracy without excessive computational cost.

The spacecraft dynamics are governed by the generalized equations of motion for
a rigid body subject to external forces and moments:
\begin{align}
    m\ddot{\mathbf{r}}                                                                               & = \mathbf{F}_{ext} \label{eq:translational_eom}  \\
    \mathbf{I}\boldsymbol{\dot{\omega}} + \boldsymbol{\omega} \times (\mathbf{I}\boldsymbol{\omega}) & = \boldsymbol{M}_{ext} \label{eq:rotational_eom}
\end{align}
where $m$ is the spacecraft mass, $\mathbf{r}$ is the position vector, $\mathbf{F}_{ext}$ represents external forces (primarily thruster forces in this study), $\mathbf{I}$ is the inertia tensor, $\boldsymbol{\omega}$ is the angular velocity vector, and $\boldsymbol{M}_{ext}$ represents external moments generated by thruster configurations.

\subsection{Spacecraft Models and Parameters}

The spacecraft models are based on the vehicles depicted in Christopher Nolan's
\textit{Interstellar}, with physical parameters derived from
\textcite{Girija2016} and supplemented by technical specifications from the
film's fan maintained wiki \parencite{fandom2025lander}.

\subsubsection{Lander 2 Configuration}

Lander 2 serves as the chaser vehicle, characterized by a compact design
optimized for atmospheric entry and orbital maneuvers. The spacecraft has a
total mass of 80,000 kg and features four faceted fuel
pods flanking its fuselage. Each pod is equipped with a set of 6 reaction
control system (RCS) hydrazine/nitrogen-tetroxide thrusters with a specific impulse of $\approx$ 330 seconds and a maximum thrust of 10000 kN per thruster.
This approximation is rather crude however and results in a very sensitive
spacecraft when it is not docked to the Endurance station. For this reason the
maximum thrust was limited to 30 - 60 kN per thruster depending on the position
and role of the thruster in the control scheme \parencite{Girija2016,muhalim2010design}.

The RCS configuration consists of 24 individual thrusters strategically
positioned to provide full 6-DOF control authority. However, the neural network
interface employs a simplified control architecture that maps six control
inputs to specific thruster groupings, as detailed in
Table~\ref{tab:control_mapping}. This was done to reduce the complexity of the
control problem and to ensure that the neural network could learn effective
maneuvers without having to also learn how to fly in a straight line first.
\begin{table}[htbp]
    \centering
    \caption{Neural Network Control Vector Mapping to Thruster Groups. Note that the directions are in the body-frame (see \autoref{chap:reference-frames}). A diagram showing the position of the thrusters can be found in Appendix~\ref{appendix:thruster-diagram}.}
    \label{tab:control_mapping}
    \begin{tabular}{clcc}
        \toprule
        \textbf{Control Input} & \textbf{Function} & \textbf{Active Thrusters} & \textbf{Max Thrust (N)} \\
        \midrule
        $u_1$                  & Translation +X    & rcs\_1, rcs\_13           & 30,000                  \\
        $u_2$                  & Translation -X    & rcs\_5, rcs\_10           & 30,000                  \\
        $u_3$                  & Rotation -Z       & rcs\_19, rcs\_23          & 60,000                  \\
        $u_4$                  & Rotation +Z       & rcs\_18, rcs\_22          & 60,000                  \\
        $u_5$                  & Translation +Y    & rcs\_9, rcs\_14           & 30,000                  \\
        $u_6$                  & Translation -Y    & rcs\_2, rcs\_6            & 30,000                  \\
        \bottomrule
    \end{tabular}
\end{table}

The spacecraft's inertia tensor reflects its elongated configuration:
\begin{equation}
    \mathbf{I}_{Lander2} = \begin{bmatrix}
        2.25 \times 10^6 & 0                & 0               \\
        0                & 2.25 \times 10^6 & 0               \\
        0                & 0                & 4.5 \times 10^6
    \end{bmatrix} \text{ kg·m}^2
\end{equation}

\subsubsection{Endurance Station Configuration}

The Endurance station represents the target spacecraft, with a much larger mass
of 1,600,000 kg and correspondingly higher moments of inertia. The station's
cylindrical configuration with rotating sections yields the inertia tensor
\parencite{Girija2016}:
\begin{equation}
    \mathbf{I}_{Endurance} = \begin{bmatrix}
        7.2 \times 10^9 & 0               & 0                \\
        0               & 7.2 \times 10^9 & 0                \\
        0               & 0               & 14.4 \times 10^9
    \end{bmatrix} \text{ kg·m}^2
\end{equation}

As the spindown maneuver was not performed in this study, this tensor ends up
not being used in the simulations.

% \subsection{Interpolation Strategy and Performance Optimization}
% Why use interpolators, speed vs accuracy trade-offs
% Major downside: no stochastic problems meaning overfitting is a risk

\subsection{State Vector and Control Architecture}

The neural network interface uses three different state vector configurations depending on problem complexity. All follow the \cite{leitner2010evolving} method of computing relative differences between chaser and target:

Basic problem (6 inputs): position differences $(x_{diff}, y_{diff})$, velocity differences $(\dot{x}_{diff}, \dot{y}_{diff})$, and angular differences $(z_{angle\_diff}, z_{rate\_diff})$.

Extended problem (8 inputs): adds travel angle (velocity-to-position alignment) and pointing angle (spacecraft orientation relative to target).

Advanced problem (10 inputs): further includes relative distance and speed magnitudes for enhanced situational awareness.

\section{Problem Formulation and Complexity Progression}
The autonomous rendezvous and docking problem was formulated as a series of
increasingly complex scenarios, progressing from basic translational maneuvers
to the restricted-in-plane problem depicted in \textit{Interstellar}. This
progression served as a software development methodology, allowing for
incremental validation and debugging of the simulation framework.

\subsection{In-Plane Problem Reduction}

Watching the \textit{Interstellar} docking sequence shows that the maneuver is
predominantly planar. The Endurance station rotates about a fixed axis
perpendicular to Cooper's approach trajectory, with the critical docking port
located on the station's perimeter. The spacecraft maintains its approach
within the plane of the station's rotation throughout the sequence, with only
small vertical motion during the final mating. The docking scenario was
constrained to in-plane motion within the X-Y coordinate system, reducing the
full 6-DOF problem to 3-DOF (two translational axes plus rotation about the
Z-axis). This reduction decreases the dimensionality of both the state space
and control space, making the problem significantly easier to solve for an
artificial neural network (ANN).

\subsection{Basic Problem: Straight-Line Translation}

In the basic problem, Lander 2 must perform a simple translational maneuver
from the origin $(0, 0)$ to reach the stationary Endurance located at $(25, 0)$
meters in the inertial frame.

\subsubsection{Problem Configuration}

The problem is set up based on inspiration from \cite{leitner2010evolving},
using a stop neuron. When the activation of this neuron
reaches a threshold, the simulation is terminated; allowing the neural network
to optimize for time. The control vector available is reduced to 3 neurons: +X
and -X, as well as the stop neuron, providing pure translational control. The
neural network architecture uses 6 input neurons based again on
\cite{leitner2010evolving}, 10 hidden neurons, and 2+1 output neurons
corresponding to the active thruster pair (the +1 being the stop
neuron).

The state vector is computed as the pose error between the chaser and target spacecraft:
\begin{equation}
    \mathbf{x}(t) = \begin{bmatrix}
        x_{target} - x_{chaser}               \\
        y_{target} - y_{chaser}               \\
        \dot{x}_{target} - \dot{x}_{chaser}   \\
        \dot{y}_{target} - \dot{y}_{chaser}   \\
        \theta_{z,target} - \theta_{z,chaser} \\
        \dot{\theta}_{z,target} - \dot{\theta}_{z,chaser}
    \end{bmatrix}
\end{equation}

To consider chaser to have docked with the target, the following constraints
must be satisfied at the end of the simulation:

\begin{align}
    \text{Position error:} \quad    & ||\mathbf{r}_{rel}|| < 0.1 \text{ m}                      \\
    \text{Velocity error:} \quad    & ||\mathbf{v}_{rel}|| < 0.1 \text{ m/s}                    \\
    \text{Orientation error:} \quad & ||\boldsymbol{\theta}_{rel}|| < \frac{\pi}{8} \text{ rad}
\end{align}

\subsubsection{Cost Function Design}
The cost function is split in two regions, one to provide selective pressure to
minimize pose error, and the second provides pressure
to minimize time. The cost function is defined as:
\begin{equation}
    C = \begin{cases}
        1 + \frac{t_{max} - t_{final}}{t_{max}}      & \text{if constraints satisfied} \\
        \frac{1}{(1 + e_{\theta})(1 + e_r)(1 + e_v)} & \text{otherwise}
    \end{cases}
\end{equation}

where $e_{\theta}$, $e_r$, and $e_v$ represent orientation, position, and
velocity errors respectively. This formulation rewards both successful docking
(fitness > 1) and penalizes constraint violations through the inverse error
product.

\subsection{Extended Problem: Diagonal Movement}

Building on the basic problem, Lander 2 must now perform a diagonal translation
from the origin $(0, 0)$ to reach the stationary Endurance at $(10, 10)$ meters,
requiring simultaneous X-Y axis coordination.

\subsubsection{Problem Configuration}

The control system was expanded first from 2 to 4 thrusters groups, +X, -X, +Rz
and -Rz providing translational rotational control control, later it was
expanded to also include +Y and -Y to provide lateral control. The neural
network architecture used the same 6-input state vector, now with between 10-40
hidden neurons, and between 4+1 and 6+1 output neurons.

\subsubsection{Cost Function Evolution}

Initially, the same dual-mode cost function from the basic problem was applied. However, this approach proved inadequate due to reward sparsity: networks would converge to local optima where they successfully minimized error along only one axis (typically X) while ignoring the other. The sparse reward structure provided insufficient guidance for coordinated multi-axis control.

To address this limitation, the cost function was redesigned using an approach inspired by \cite{park2017nonlinear}, where an ANN was trained on a similar problem.
The new cost function incorporated shaping rewards to provide continuous feedback on progress towards the goal, rather than a binary success/failure signal. The revised cost function is defined as:

\begin{equation}
    J = (\mathbf{x}_{final} - \mathbf{x}_{target})^T P (\mathbf{x}_{final} - \mathbf{x}_{target}) + \int_{0}^{t_{final}} (\mathbf{x}(t) - \mathbf{x}_{target})^T Q (\mathbf{x}(t) - \mathbf{x}_{target}) dt
\end{equation}
Where $P$ should be the positive-definite solution of the discrete Riccati equation for the linearized system, and $Q$ is a positive-definite matrix that weights the state error over time. Matrix $P$ was empirically set to $3*Q$ for simplicity, while $Q$ was set using \cite{bryson2018applied}'s rule to normalize the state errors using these weighting factors. The tolerances for position, velocity, and orientation errors remained the same as in the basic problem's docking tolerances, and the tolerance for angular velocity was set to $0.01$ rad/s to push hard for minimizing rotational difference.
\\
To make the cost function even more dense, shaping rewards as well as penalties were added to encourage intermediate milestones and punish 'cheating' such as spinning in place to reduce angular velocity error without translating. These shaping rewards included:
\begin{itemize}
    \item A reward for reducing the angle between the velocity vector and the line-of-sight
          to the target:\\ $R_{align} = w_{align} \cdot \max(0, \hat{\mathbf{v}} \cdot \hat{\mathbf{r}}_{los}) \cdot dt$ \\where $w_{align} = 10.0$
    \item A reward for reducing the distance to the target: \\$R_{progress} = w_{progress} \cdot (d_{prev} - d_{curr})$ \\where $w_{progress} = 20.0$
    \item A reward for initiating rotation: \\$R_{spin\_init} = w_{spin\_init} \cdot dt$ where $w_{spin\_init} = 5.0$ if $||\boldsymbol{\omega}|| > 10^{-3}$
    \item A reward for reducing angular velocity error: \\$R_{spin\_match} = w_{spin\_match} \cdot \delta\omega \cdot dt$ \\where $w_{spin\_match} = 15.0$ if $\delta\omega > 0$
    \item A penalty for early termination: \\$P_{early} = w_{early} \cdot \frac{t_{planned} - t_{actual}}{t_{planned}}$ \\where $w_{early} = 5 \times 10^7$ if $t_{actual} < 0.9 t_{planned}$
    \item A penalty for lingering near start position: \\$P_{linger} = w_{linger} \cdot dt$ where $w_{linger} = 5 \times 10^4$ \\if $t > t_{start} + t_{grace}$ and $d_{current} > 0.8 \cdot d_{start}$
    \item A penalty for excessive final angular velocity: \\$P_{spin\_excess} = w_{linger} \cdot (||\boldsymbol{\omega}_{rel}|| - \omega_{tol})$ \\where $w_{linger} = 5 \times 10^4$ if $||\boldsymbol{\omega}_{rel}|| > 3\omega_{tol}$
\end{itemize}



\subsection{Advanced Problem: Rotation and Combined Maneuvers}

The final and most challenging problem configuration represented the full
restricted in-plane docking scenario inspired by the \textit{Interstellar}
sequence. This formulation required Lander 2 to perform first the translation
to the docking port, and then a spinup maneuver to match the angular phase and
velocity of Endurance (although the phase constraint was dropped later to
reduce complexity). The cost function remained the same as in the extended
problem.

\begin{align}
    ||\boldsymbol{\Omega}_{rel}|| < 0.1 \text{ rad/s} \\
    ||\boldsymbol{\omega}_{rel}|| < 0.1 \text{ rad/s}
\end{align}

\section{Algorithm Evolution and Selection}
\subsection{Parameter Exploration and Sensitivity Analysis}

Like mentioned in the theory section, an island model requires experimentation to determine if it is beneficial for the specific problem. SGA was used with an island model to investigate the effects of migration frequency and population distribution.

The parameters investigated were migration frequency (every generation vs. multiple generations between migrations) and population size versus generation count trade-offs. Each run was configured to use the same amount of evaluations (3600) and were compared by average and maximum fitness using the basic problem. This exploration revealed that large gene pools with frequent migration converged quickly but achieved lower average fitness, while allowing more generations between migrations yielded better diversity maintenance and superior final average fitness. This highlighted the importance of diversity preservation.

\subsection{Initial Approaches: SGA and Island Model}

Simple Genetic Algorithm successfully solved the basic straight-line translation problem, with the island model variant achieving faster convergence through improved diversity maintenance. However, both approaches failed on the diagonal movement problem, likely due to the combination of sparse reward signals from the original cost function and insufficient diversity preservation mechanisms to escape local optima in the higher-dimensional search space.

\subsection{Transition to NEAT}

NEAT was implemented due to uncertainty about optimal network topology. Varying neuron count with SGA was ineffective, and expanding the search space to include layer variations would have been computationally expensive. NEAT's speciation mechanism addressed the diversity issues discovered in SGA experiments.

Initially, NEAT struggled with the diagonal problem using the original cost function. However, after implementing the improved cost function inspired by Park et al., NEAT successfully solved the diagonal movement problem. Unfortunately, NEAT was unable to solve the full rotating target problem within the available time constraints. 

NEAT was however able to solve an easier version of the rotating problem, where the target was positioned only 3 meters away and rotating at a reduced angular velocity. This was achieved by taking the population that had solved the diagonal problem and retraining them a pure phase matching problem, where the chaser had to match the angular velocity and phase of a stationary target placed placed on top of the chaser. Within a few generations, several networks were able to solve this problem. These networks were then used as the starting population for the easier problem and were able to solve it. This demonstrated that NEAT is capable of learning the necessary behaviors, but struggles to learn them all at once.


